\emph{Mixed-multiplicity spectra of block-equalized coincident copies.} In the rank-one edge-datum family of the third paragraph above, group \(q=q_1+\cdots+q_s\) coincident copies into \(s\) factor blocks and take the pencil data of the second paragraph above block-scalar, \(H=\operatorname{Im}C_0=\bigoplus_f h_f\mathbf 1_{q_f}\) and \(K=Z_0-C_0^\dagger C_0=\bigoplus_f n_f\mathbf 1_{q_f}\) with \(n_f>0\) - the per-factor relaxation of exact norm equalization, obtained by taking the momenta equal within each factor block and unequal across blocks. Then for every \(q\) and every \(N\ge2\) the pencil characteristic polynomial factors exactly as \(\det(\lambda\mathbf 1_{2q}-T)=\prod_f(\lambda^{2}+4\lambda h_f-4n_f)^{q_f}\), so the nonzero spectrum of \(\widetilde Q\) is the multiset union over \(f\) of \(q_f\) copies of the Row length pair at \((h_f,n_f)\), each block carrying its own involution \(\lambda\mapsto-4n_f/\lambda\), its own row length \(2n_f^{1/2}\) and its own row-pair splitting \(4(h_f^{2}+n_f)^{1/2}\) - the general-\(q\) equalized spectrum of the second paragraph above, extended from uniform depth to per-factor depths (exact; symbolic at \(q=2\) for the partitions \((1,1)\) and \((2)\) and at \(q=3\) for the partition \((2,1)\); numerical to \(1.5\times10^{-14}\) through \(q=5\) and \(N\le8\) on edge blocks constructed explicitly as \(R=U(A+\mathrm iH)+V\,K^{1/2}\) with \(V\) an isometry into the orthogonal complement of the column space of \(U\), so that every datum of the class is realized at every \(q\): a data-level construction, not a statement that generic ensembles produce it, and it realizes the non-uniform per-factor depths that the copy-invariant-momenta class of the fourth paragraph above excludes, whose hypothesis of invariance across the whole copy space is here relaxed to invariance within each factor block). Within the class, the exact \(q=2\) criterion of the third paragraph above is recovered on the positive-Gram locus: writing the quartic as \(\lambda^{4}+a\lambda^{3}+b\lambda^{2}+c\lambda+d\), the two conditions \(8c=a(4b-a^{2})\) and \(64d=(4b-a^{2})^{2}\), which are necessary and sufficient for it to be a perfect square of a real monic quadratic, hold jointly if and only if \(h_1=h_2\) and \(n_1=n_2\) - the second condition factors as \(1024n_1n_2-256\bigl((h_1-h_2)^{2}+n_1+n_2\bigr)^{2}\), and \((h_1-h_2)^{2}+n_1+n_2\ge2(n_1n_2)^{1/2}\) with equality only at coincidence; the off-locus algebraic branch \(n_f=-h_f^{2}\) lies outside \(K\succeq0\), the double-pair sign is automatic on positive block-diagonal data since \(4b-a^{2}=-16((h_1-h_2)^{2}+n_1+n_2)\), the mismatched double-root example \((3,7)\), \((5,11)\) of the second paragraph above fails both conditions with values \(2560\) and \(-45056\), and none of \(2000\) random general \((H,K)\) samples meets them (exact factorizations and rational arithmetic; numerical sampling). At \(h_f=0\) the ratio of the mixed-multiplicity moment \(\sum_j e_j^{2p}\) of the nonzero spectrum to the multiplicity-\(1\) moment equals \(1+\sum_f(q_f-1)n_f^{p}/\sum_f n_f^{p}\), the signature-excess formula of the fourth paragraph above with \(n_f\) in place of \(w_f\): \(9/17\) at \(p=1\), \(81/145\) at \(p=2\) at \((q_1,q_2)=(2,1)\) and \((n_1,n_2)=(1/4,\,2/9)\), and \(1\) at equal depths (exact rational arithmetic); at \(h_f\ne0\) the same ratio holds with the pair \(2p\)-moments \(8(2h_f^{2}+n_f)\) and \(32(8h_f^{4}+8h_f^{2}n_f+n_f^{2})\) as weights, verified from the full spectra to \(5\times10^{-15}\), so in this class the stacked-versus-multiplicity-\(1\) moment excess is carried by the charge spectrum itself - an identity of moment signatures, not a selection statement. The factorization is exact algebra at every \(q\) and \(N\ge2\), while the reading of spectral multiplicity as \(q_f\) equal rows is asserted, as in the third and fourth paragraphs above, only at fixed depths and leading semiclassical order, the regime of depths growing with \(N\) being untreated. The relative-gap diagnostic of the third paragraph above vanishes if and only if all per-factor pairs \((h_f,n_f)\) coincide - for sorted nonzero eigenvalues the multiplicity of each extreme value is strictly below \(q\) at genuine mismatch, so one of its two defining gaps is then strictly positive - and it is at least \(9.5\times10^{-2}\) on the sampled mixed grid while it vanishes at coincidence, so non-uniform depths leave a nonvanishing relative gap in the charge spectrum; the fourth-moment signature excess stays nonnegative and vanishes exactly when all nonzero eigenvalues have equal modulus (numerical, on the full charge). These are exact spectrum-structure, criterion and moment-identity statements for the block-scalar subclass of the rank-one edge-datum family at fixed \(q\) and arbitrary \(N\), numerical where marked so; they are statements about pencil and charge spectra, not about any state or ensemble, they leave the two-sided \(\pm\) reading, the typicality question and the general-\((H,K)\) criterion outside the block-scalar subclass exactly as the two paragraphs above, they treat per-factor depths only in the fixed-depth regime of the selection-rule reading, and they close nothing in Problem~G.
